\chapter{Organisation}

\section{Organisational chart}

The chart was drawn from the leadership page of AstraZeneca \cite{LEAD}. It shows the Board of Directors, the Chief Executive Officer (CEO) and the Senior Executive Team (SET) in Figure~\ref{fig:unidades}. The source for each box is given in Table~\ref{tab:fontes-organograma}.

\begin{figure}[H]\centering
\resizebox{\linewidth}{!}{%
\begin{tikzpicture}
  \node[orgtop] (b) at (0,0) {\nd{Board of Directors}{14 directors · 5 committees}};
  \node[org=2.45cm] (d0) at (-6.75,-1.30) {\nd{M. Demaré}{Chair}};
  \node[org=2.45cm] (d1) at (-4.05,-1.30) {\nd{P. Broadley}{Senior Independent}};
  \node[org=2.45cm] (d2) at (-1.35,-1.30) {\nd{E. Ashley}{}};
  \node[org=2.45cm] (d3) at (1.35,-1.30) {\nd{B. Conix}{}};
  \node[org=2.45cm] (d4) at (4.05,-1.30) {\nd{K. Knudsen}{}};
  \node[org=2.45cm] (d5) at (6.75,-1.30) {\nd{D. Layfield}{}};
  \node[org=2.45cm] (d6) at (-6.75,-2.6) {\nd{A. Manz}{}};
  \node[org=2.45cm] (d7) at (-4.05,-2.6) {\nd{S. McCoy}{}};
  \node[org=2.45cm] (d8) at (-1.35,-2.6) {\nd{T. Mok}{}};
  \node[org=2.45cm] (d9) at (1.35,-2.6) {\nd{M. Wallenberg}{}};
  \node[org=2.45cm] (d10) at (4.05,-2.6) {\nd{R. Haas}{}};
  \node[org=2.45cm] (d11) at (6.75,-2.6) {\nd{W. L. Martello}{}};
  \node[org=3.2cm] (ceo) at (0,-3.5) {\nd{Pascal Soriot}{Chief Executive Officer}};
  \node[org=3.2cm] (set) at (0,-4.7) {\nd{Senior Executive Team}{10 members}};
  \node[org=1.80cm] (c0) at (-7.88,-6.2) {\nd{Aradhana Sarin}{CFO}};
  \node[org=2.10cm] (c1) at (-5.41,-6.2) {\nd{Oncology}{}};
  \node[org=2.10cm, below=6pt of c1] (c10) {\nd{D. Fredrickson}{EVP, Oncology Haematology Business Unit}};
  \draw[orglink] (c1.south) -- (c10.north);
  \node[org=2.10cm, below=6pt of c10] (c11) {\nd{S. Galbraith}{EVP, Oncology Haematology R\&D}};
  \node[org=2.55cm] (c2) at (-2.57,-6.2) {\nd{BioPharmaceuticals}{}};
  \node[org=2.55cm, below=6pt of c2] (c20) {\nd{S. Barr}{EVP, BioPharma R\&D}};
  \draw[orglink] (c2.south) -- (c20.north);
  \node[org=2.55cm, below=6pt of c20] (c21) {\nd{R. Dobber}{EVP, BioPharma BU}};
  \node[org=2.10cm] (c3) at (0.27,-6.2) {\nd{Rare Disease}{}};
  \node[org=2.10cm, below=6pt of c3] (c30) {\nd{M. Dunoyer}{CEO Alexion \& Chief Strategy Officer}};
  \draw[orglink] (c3.south) -- (c30.north);
  \node[org=1.90cm] (c4) at (2.79,-6.2) {\nd{I. Reic}{EVP, International}};
  \node[org=2.00cm] (c5) at (5.26,-6.2) {\nd{P. Cheng}{EVP Global Operations, IT \& CSO}};
  \node[org=2.00cm] (c6) at (7.78,-6.2) {\nd{J. Pott}{CHRO, Compliance \& General Counsel}};
  \coordinate (bus) at (0,-5.85);
  \draw[orglink] (b.south) -- (ceo.north);
  \draw[orglink] (ceo.south) -- (set.north);
  \draw[orglink] (set.south) -- (bus);
  \draw[orglink] (c0.north |- bus) -- (c6.north |- bus);
  \draw[orglink] (c0.north |- bus) -- (c0.north);
  \draw[orglink] (c1.north |- bus) -- (c1.north);
  \draw[orglink] (c2.north |- bus) -- (c2.north);
  \draw[orglink] (c3.north |- bus) -- (c3.north);
  \draw[orglink] (c4.north |- bus) -- (c4.north);
  \draw[orglink] (c5.north |- bus) -- (c5.north);
  \draw[orglink] (c6.north |- bus) -- (c6.north);
\end{tikzpicture}}

\caption{AstraZeneca organisational chart, redrawn for this report. Source: authors' elaboration, based on \cite{LEAD}. Every box is stated in the AR \cite[pp.~67--70]{AR}, except the Board change made after it was filed \cite{AGM26,MART}; the source for each box is given in Table~\ref{tab:fontes-organograma}.}
\label{fig:unidades}
\end{figure}

\begin{xltabular}{\linewidth}{@{}>{\raggedright\arraybackslash}p{3.4cm}Y@{}}
\caption{Sources for each box of Figure~\ref{fig:unidades}.}\label{tab:fontes-organograma}\\*
\toprule
\textbf{Box} & \textbf{Source and notes} \\*
\midrule
\endfirsthead
\toprule
\textbf{Box} & \textbf{Source and notes} \\*
\midrule
\endhead
\bottomrule
\endfoot
Board of Directors & The Board has 14 directors and five committees (Audit, Remuneration, Nomination and Governance, Science, Sustainability). Matters not reserved to the Board are delegated to the CEO or to one of the committees \cite[pp.~67--68]{AR}. Since the AR, Nazneen Rahman retired at the 2026 AGM \cite{AGM26} and Wan Ling Martello joined on 1 September 2026 \cite{MART}, so the Board still has 14 directors. \\
Michel Demaré & Non-Executive Chair of the Board \cite[p.~68]{AR}. \\
Philip Broadley & Senior Independent Non-Executive Director \cite[p.~68]{AR}. \\
Euan Ashley, Birgit Conix, Karen Knudsen, Diana Layfield, Anna Manz, Sheri McCoy, Tony Mok, Marcus Wallenberg & Listed as Non-Executive Directors in the AR \cite[pp.~68--69]{AR}. \\
Rene Haas & Non-Executive Director since 1 January 2025 \cite[pp.~68--69]{AR}. Re-elected at the 2026 AGM with 97.97\% of votes cast \cite{AGM26}. \\
Wan Ling Martello & Non-Executive Director since 1 September 2026, joining the Audit Committee. Appointed after the AR (FY2025) was filed, which is why she is not listed in it \cite{MART}. \\
Pascal Soriot, CEO & Executive Director and Chief Executive Officer; member of the SET \cite[pp.~68, 70]{AR}. \\
Senior Executive Team (SET) & The SET is the body through which the CEO exercises the authority delegated by the Board. It has 10 members, five of them women \cite[pp.~39, 70]{AR}. \\
Aradhana Sarin & Executive Director and Chief Financial Officer; a SET member \cite[pp.~68, 70]{AR}. \\
Oncology & One of the three R\&D organisations named in the AR \cite[p.~26]{AR}. Both EVPs below carry ``Oncology Haematology'' in their own AR title \cite[p.~70]{AR}. \\
David Fredrickson & Executive Vice President, Oncology Haematology Business Unit \cite[p.~70]{AR}. \\
Susan Galbraith & Executive Vice President, Oncology Haematology R\&D \cite[p.~70]{AR}. \\
BioPharmaceuticals & One of the three R\&D organisations named in the AR \cite[p.~26]{AR}. Both EVPs below carry ``BioPharmaceuticals'' in their own AR title \cite[p.~70]{AR}. \\
Sharon Barr & EVP, BioPharmaceuticals R\&D \cite[p.~70]{AR}. \\
Ruud Dobber & Executive Vice President, BioPharmaceuticals Business Unit, responsible for CVRM, R\&I and Vaccines \& Immune Therapies together \cite[p.~70]{AR}. \\
Rare Disease & One of the three R\&D organisations named in the AR \cite[p.~26]{AR}. Alexion, acquired in 2021, is AstraZeneca's rare disease business \cite[Item~4.A]{F20}; it is run within the SET by Marc Dunoyer \cite[p.~70]{AR}. \\
Marc Dunoyer & CEO of Alexion and Chief Strategy Officer of AstraZeneca, a SET member. Two official posts held by one person \cite[p.~70]{AR}. \\
Iskra Reic & EVP, International: China, Asian and Eurasian markets, Middle East \& Africa, Latin America, Australia and NZ. International is one of the AR's three commercial regions, and the only one with its own named SET member; the US and Europe-Canada regions have no separate SET-level head in the AR \cite[pp.~32, 70]{AR}. \\
Pam Cheng & EVP, Global Operations, IT and Chief Sustainability Officer \cite[p.~70]{AR}. \\
Jeff Pott & Chief Human Resources Officer, Chief Compliance Officer and General Counsel (the AR also lists IP) \cite[p.~70]{AR}. \\
\end{xltabular}

Apart from the Board change noted at the end of this paragraph, every box corresponds to a fact stated in the AR: the chain Board $\rightarrow$ CEO $\rightarrow$ SET \cite[pp.~67, 70]{AR}; the ten SET members and their titles \cite[p.~70]{AR}; and the three therapy-area R\&D organisations stated in the AR \cite[p.~26]{AR}. Oncology and BioPharmaceuticals are confirmed further by the words of their own EVPs' AR titles \cite[p.~70]{AR}; Rare Disease is confirmed instead through Alexion, which the Form 20-F describes as AstraZeneca's rare disease business \cite[Item~4.A]{F20} and which Marc Dunoyer runs as a SET member \cite[p.~70]{AR}. Iskra Reic's EVP, International, is one of the AR's three commercial regions, and the only one with a named SET member: the AR does not give the US or Europe-Canada regions a separate SET-level head \cite[pp.~32, 70]{AR}. No box groups unrelated SET members under a heading the AR does not use. The Board is shown as it stands in October 2026: Nazneen Rahman, listed in the AR, retired at the 2026 AGM \cite{AGM26}, and Wan Ling Martello was appointed on 1 September 2026 \cite{MART}.

\section{Analysis of the structure}

The criteria below come from the course material on the organising function \cite{ORG}. Each statement is limited to what the AR, the chart or the course material support, and each analysis ends with what cannot be concluded from the sources.

\subsection{Type of departmentalisation}

Departmentalisation means grouping activities into logical units to carry out the strategy \cite[p.~14]{ORG}. The course material distinguishes functional grouping \cite[p.~15]{ORG}, divisional grouping by product, area, customer or project \cite[pp.~16--19]{ORG}, multiple (mixed) grouping when large organisations cannot rely on one criterion \cite[p.~20]{ORG}, and matrix structures, which combine two criteria and have dual authority, with employees reporting to two managers \cite[pp.~21--22]{ORG}.

AstraZeneca uses several criteria at the same time:

\begin{itemize}
\item \textbf{By therapy area (product):} three R\&D organisations, Oncology, BioPharmaceuticals and Rare Disease \cite[p.~26]{AR}, with a separate EVP for R\&D and for the business unit in Oncology and in BioPharmaceuticals \cite[p.~70]{AR}.
\item \textbf{By geography:} three commercial regions (US, Europe-Canada, International), with Japan reporting separately \cite[p.~32]{AR}; only International has its own named SET member, Iskra Reic \cite[p.~70]{AR}.
\item \textbf{By function:} Operations (manufacturing and supply, IT), finance, and human resources, compliance and legal are led by separate SET members \cite[p.~70]{AR}. AstraZeneca's own careers site independently names Operations as one of its business areas, covering ``the development, manufacturing, testing and delivery of our medicines'' \cite{CAREERS}.
\end{itemize}

\begin{analysis}
The structure is best described as \textbf{multiple (mixed) departmentalisation}, combining the functional, therapy-area and geographic criteria \cite[p.~20]{ORG}.
\end{analysis}

\begin{sourceslimit}
It cannot be called a matrix. A matrix requires dual authority \cite[pp.~21--22]{ORG}, and neither the chart nor the AR documents anyone reporting to two managers: Ruud Dobber's single role as EVP, BioPharmaceuticals Business Unit covers CVRM, R\&I and Vaccines \& Immune Therapies together \cite[p.~70]{AR}, which is one person covering several areas, not dual reporting. The AR also says the Group has one reportable segment and that discovery, development, manufacturing and commercialisation are managed globally on a highly integrated basis \cite[p.~147]{AR}.
\end{sourceslimit}

\subsection{Capacity to adapt to change}

The course material links the organising function to change management \cite[pp.~4--5]{ORG}, says that structures with few levels and delegation respond faster than those with many levels and centralisation \cite[p.~7]{ORG}, and that functional structures are rigid and hard to restructure \cite[p.~16]{ORG}.

The AR documents deliberate, long-running change:

\begin{itemize}
\item After acquiring Alexion in 2021, the Group launched the Post Alexion Acquisition Group Review, a global restructuring programme to integrate systems, structure and processes. It cost \$232m in 2025 (of \$237m of total restructuring costs) and \$3,414m cumulatively \cite[p.~142]{AR}, with expected total costs of \$4.0bn \cite[p.~59]{AR}.
\item ``Enable an agile organisation'' is one of the strategic focus areas, alongside ``Transform R\&D ways of working'' \cite[p.~11]{AR}.
\item In December 2025 the Group created a dedicated Enterprise AI unit to accelerate Ambition 2030 \cite[p.~36]{AR}.
\item Marc Dunoyer leads Alexion within the SET \cite[p.~70]{AR}. Before her current role, the CFO, Aradhana Sarin, was Alexion's Chief Financial Officer \cite{SARIN21}, and Sharon Barr was Alexion's SVP, Head of Research and Product Development \cite{BARR23}; both are confirmed by AstraZeneca's own appointment announcements, not only by the AR.
\end{itemize}

\begin{analysis}
The Group has shown it can absorb a large acquisition and reorganise on purpose, and it names agility as a goal. This points to a \textbf{reasonable capacity to adapt}.
\end{analysis}

\begin{sourceslimit}
The AR gives no decision times, so the \textit{speed} of response cannot be assessed. Career moves of individual executives (for example Iskra Reic) and a ``2024 reorganisation'' are not in the AR, so they are not used.
\end{sourceslimit}

\subsection{Formal and informal organisation}

The formal organisation is designed by management, has rules and procedures, and can be drawn in an organisational chart. The informal organisation emerges from social interaction, is not and cannot be in the chart, and can speed communication or resist change \cite[pp.~12--13]{ORG}.

\textbf{Formal:} the Board keeps certain matters reserved and delegates the rest to the CEO or to its five committees \cite[p.~67]{AR}; the SET is the body through which the CEO exercises delegated authority \cite[p.~70]{AR}; there is a Code of Ethics and a global compliance programme \cite[p.~34]{AR}; and cross-functional sub-committees and product committees allocate resources \cite[p.~147]{AR}. Marc Dunoyer holding two posts, CEO of Alexion and Chief Strategy Officer, is also \textit{formal}, because both are official positions \cite[p.~70]{AR}.

\textbf{Closest documented link to informal organisation:} the Employee Resource Groups, which are voluntary, employee-led groups open to all employees \cite[p.~39]{AR}, and the Pulse survey of employees \cite[p.~39]{AR}.

\begin{sourceslimit}
No document describes the real informal organisation (networks, communities of practice, informal leaders). By definition the chart does not show it \cite[p.~12]{ORG}. The report therefore only explains the concept and the organised channels that exist.
\end{sourceslimit}

\subsection{Flexibility or rigidity}

According to the course material, large organisations tend to be more specialised, centralised and formalised \cite[p.~11]{ORG}, and rigidity grows with formalisation \cite[p.~9]{ORG} and centralisation \cite[p.~7]{ORG}.

\begin{itemize}
\item All significant operating decisions are taken at SET level, and resources are allocated Group-wide by product committees \cite[p.~147]{AR}. This indicates \textbf{centralisation} at the top.
\item Governance is documented in detail: reserved matters, delegations, committees and their roles \cite[pp.~67, 71]{AR}.
\item Commercial presence is organised through wholly-owned local marketing companies in most markets \cite[p.~32]{AR}, and the Group states a focus on an agile organisation \cite[p.~11]{AR}.
\end{itemize}

\begin{analysis}
The evidence points to a structure that is \textbf{formalised and centralised at the SET}, with specialisation by therapy area and region, and with agility as a declared goal. It leans towards rigidity in decision-making and towards flexibility only in the stated intention.
\end{analysis}

\begin{sourceslimit}
The AR does not describe autonomy for units below the SET, nor regional pricing autonomy. The idea that Alexion works as an ``autonomous company'' is not supported: the AR says the Group has a single reportable segment and that functional areas are not managed separately \cite[p.~147]{AR}.
\end{sourceslimit}

\subsection{Number of hierarchical levels}

The course material relates levels to span of control: the wider the span, the fewer the levels and the flatter the structure; narrower spans give more levels and a tall structure \cite[pp.~8--9]{ORG}. It also names three management levels: strategic, tactical and operational \cite[p.~7]{ORG}.

\begin{table}[H]\centering
\caption{Levels counted in Figure~\ref{fig:unidades}, including the Board as level~1.}
\label{tab:niveis}
\begin{tabularx}{\linewidth}{@{}>{\raggedright\arraybackslash}p{4.2cm}cY@{}}
\toprule
\textbf{Branch} & \textbf{Levels} & \textbf{Chain} \\
\midrule
CFO, International, Global Operations and HR/Legal & 4 & Board $\rightarrow$ CEO $\rightarrow$ SET $\rightarrow$ SET member (3 below the Board) \\
\addlinespace
R\&D (Oncology, BioPharmaceuticals, Rare Disease) & 5 & Board $\rightarrow$ CEO $\rightarrow$ SET $\rightarrow$ R\&D organisation $\rightarrow$ EVP (4 below the Board), because the AR itself names these three organisations \cite[p.~26]{AR} \\
\bottomrule
\end{tabularx}
\end{table}
\begin{sourceslimit}
The count refers to the \textbf{chart analysed}, which covers only the top of the company. The AR does not disclose levels below the SET nor spans of control, so it cannot be concluded whether the whole structure is tall or flat. Other level counts found in secondary sources (6 to 9 levels) are not supported by the AR and are not used.
\end{sourceslimit}

\subsection{Line and staff positions}

A \textbf{line} manager has decision power and authority over the workers who report to him or her, within the chain of command. A \textbf{staff} unit or person supports, advises or assists line managers. In the course material's example, directors of marketing, finance and production are line positions, while the board's secretariat and legal advice are staff \cite[p.~8]{ORG}.

\begin{xltabular}{\linewidth}{@{}YYY@{}}
\caption{Line and staff positions in the Senior Executive Team. Classification: authors' analysis using \cite[p.~8]{ORG}.}\label{tab:linha-staff}\\*
\toprule
\textbf{Position} & \textbf{Responsibility according to the AR} & \textbf{Classification} \\*
\midrule
\endfirsthead
\toprule
\textbf{Position} & \textbf{Responsibility according to the AR} & \textbf{Classification} \\*
\midrule
\endhead
\bottomrule
\endfoot
EVPs of R\&D and business units: Fredrickson, Galbraith, Barr, Dobber & Run the R\&D and commercial performance of their therapy areas \cite[p.~70]{AR} & \textbf{Line} \\
Iskra Reic, EVP International & Strategy and growth in the International region \cite[p.~70]{AR} & \textbf{Line} \\
Pam Cheng, EVP Global Operations, IT and CSO & Operations, IT and sustainability \cite[p.~70]{AR} & \textbf{Mixed} manufacturing is line; IT and sustainability support other areas \\
Jeff Pott, CHRO, Compliance and General Counsel & HR, compliance, legal and IP \cite[p.~70]{AR} & \textbf{Staff} legal advice is staff in the course example \\
Aradhana Sarin (CFO); Marc Dunoyer (Alexion, Chief Strategy Officer) & Finance; Alexion and strategy \cite[p.~70]{AR} & \textbf{Not classified} the AR does not say who reports to them \\
\end{xltabular}

\section{Proposed changes to the structure}

The proposals below each rest on a fact stated in the AR and on the trade-offs described in the course material.

\subsection{Proposal 1 · Stronger coordination between R\&D and the business unit in each therapy area}

\textbf{Basis.} In Oncology and in BioPharmaceuticals, R\&D and the business unit are led by different SET members \cite[p.~70]{AR}. The AR states that cross-functional sub-committees and product committees already exist \cite[p.~147]{AR}.

\textbf{Proposal.} Make the coordination at therapy-area level permanent and visible, with one joint forum per area where R\&D and the business unit share priorities and resources.

\textbf{Advantage.} Faster decisions and fewer gaps between research and commercial plans. \textbf{Risk.} More meetings and slower decisions if the forum has no clear authority. The course material warns that, with product-based units, each unit's interests can come before the organisation's and units may compete instead of cooperating \cite[p.~17]{ORG}; a joint forum is meant to counter that.

\subsection{Proposal 2 · More local autonomy in access and pricing}

\textbf{Basis.} Pricing pressure from payers is a main commercial concern in several regions \cite[p.~32]{AR}. Today the Group applies a global tiered pricing policy with four tiers based on gross national income \cite[p.~41]{AR}, and operating decisions are taken at SET level \cite[p.~147]{AR}.

\textbf{Proposal.} Give regional commercial teams more room to adapt access and pricing, within limits set centrally by the SET.

\textbf{Advantage.} Delegation brings decisions closer to the market and speeds the response \cite[p.~7]{ORG}, and when units are far apart it is more suitable to assign management to a manager located in the region \cite[p.~17]{ORG}. \textbf{Risk.} Less consistency between countries; the limits must be clear. The AR does not report how much autonomy regions have today, so this proposal needs that information to be tested.
